\chapter*{Περίληψη}
\addcontentsline{toc}{chapter}{Περίληψη}
Τα \en{NFTs} (\en{Non-Fungible Tokens}) καθιέρωσαν την έννοια του μοναδικού ψηφιακού
περιουσιακού στοιχείου σε δημόσια \en{blockchains}, όμως η αγορά τους παραμένει
οργανωμένη σχεδόν αποκλειστικά γύρω από πωλήσεις έναντι κρυπτονομίσματος: σταθερές τιμές,
δημοπρασίες και προσφορές. Η απευθείας ανταλλαγή \en{NFT} με \en{NFT}, ιδίως όταν τα δύο
αντικείμενα ανήκουν σε διαφορετικές συλλογές, δεν εξυπηρετείται από τις υπάρχουσες
πλατφόρμες και, όπου επιχειρείται, προϋποθέτει εμπιστοσύνη προς τον αντισυμβαλλόμενο ή
προς έναν μεσάζοντα.

Η παρούσα διπλωματική εργασία σχεδιάζει και υλοποιεί ένα \en{marketplace} ανταλλαγών
(\en{barter}) στο \en{Sui blockchain}, στο οποίο ένας χρήστης υποβάλλει προσφορά
(\en{bid}) που αποτελείται από ένα ή περισσότερα \en{NFTs} οποιασδήποτε συλλογής και,
προαιρετικά, ποσό \en{SUI}, με αντάλλαγμα ένα συγκεκριμένο \en{NFT} άλλου χρήστη. Τα
προσφερόμενα στοιχεία δεσμεύονται (\en{escrow}) μέσα σε ένα κοινόχρηστο αντικείμενο
\en{Bid} από τη στιγμή της υποβολής, ώστε η αποδοχή από τον κάτοχο του ζητούμενου
\en{NFT} να ολοκληρώνει την ανταλλαγή ατομικά και χωρίς θεματοφύλακα. Το έξυπνο συμβόλαιο
είναι γραμμένο στη γλώσσα \en{Move} και αξιοποιεί το αντικειμενοκεντρικό μοντέλο του
\en{Sui}.

Το ανταλλακτικό εμπόριο υπήρξε ιστορικά η πρώτη μορφή οικονομικής συναλλαγής, πολύ πριν
το χρήμα καθιερωθεί ως κοινό μέτρο αξίας, και δεν έπαψε ποτέ να έχει θέση εκεί όπου τα
αγαθά είναι μοναδικά και η αξία τους υποκειμενική. Η ψηφιοποίηση πραγματικών αγαθών ως
\en{NFTs} έχει ήδη περάσει από το πειραματικό στο θεσμικό στάδιο, με χαρακτηριστικό
παράδειγμα την ψηφιοποίηση (\en{tokenization}) τίτλων ιδιοκτησίας ακινήτων σε \en{blockchain}
από το Κτηματολόγιο του Ντουμπάι. Σε αυτό το
πλαίσιο, η ύπαρξη ενός αξιόπιστου μηχανισμού απευθείας ανταλλαγής, που δεν επιβάλλει τη
ρευστοποίηση κάθε αγαθού σε κρυπτονόμισμα ως ενδιάμεσο βήμα, παύει να αποτελεί
συμπλήρωμα των υπαρχουσών αγορών και γίνεται αναγκαία υποδομή. Η παρούσα εργασία
προτείνει έναν τέτοιο μηχανισμό, στον οποίο η ανταλλαγή εκτελείται ατομικά από το ίδιο
το συμβόλαιο, ώστε κανένα από τα δύο μέρη να μη χρειάζεται να εμπιστευθεί το άλλο.

\vspace{1cm}
\noindent\textbf{Λέξεις κλειδιά:} \en{blockchain, Sui, Move, NFT, marketplace, barter, bid,
escrow, atomic swap, smart contract, dApp, cross-collection}
\clearpage
